\chapter{Evaluation}
    - Methodologie
        - welche benchmarks, welche messungen, etc.
        Benchmark: Collatz-Reihenberechnung von 1 bis 1000, LED wechselt Phase nach jedem erfolgreichen Durchlauf der 1000 Reihen
        Messung in der Simulation: In der Simulation wird jeweils ausgegeben, wenn LED umschaltet. Ein Python-Skript liest diese Ausgabe und berechnet die Zeit zwischen je zwei Ausgaben.
        Messung im FPGA: Kamera mit 120 Frames pro Sekunde filmt Basys 3-Board und eine Stoppuhr. Es werden für die Dauer einer LED-Phase (an/aus) die Frames gezählt und mit der abgelaufenen Zeit auf der Stoppuhr verglichen.

    - Präsentation und Interpretation der Ergebnisse
        3.503 Sekunden im Mittel bei der Simulation (neuer Wert ohne Prints) bei 100 Messungen. Minimum: 3.428, Maximum: 3.873
        Auf FPGA 0.916 Sekunden im Mittel (Stand: 25 Messungen), Minimum: 0,83 Sekunden, Maximum: 0,99 Sekunden \todo{Frames zählen}
    - Bewertung
        ca. 3,8-fache Beschleunigung auf "echtem" FPGA --> FPGA ist vermutlich immer schneller, für ein stärkeres Speedup werden jedoch ggf. teurere Boards benötigt